% HIKET -- NEXT SESSION: the level/trend discrepancy and the hypotheses for it.
%
% Companion to M&M_parameterization_working_document.pdf. That document states the
% parameterisation assumptions and their evidence; this one states the problem that
% remains after them, the candidate explanations, and what would test each.
\documentclass[11pt]{article}
\usepackage[utf8]{inputenc}\usepackage[T1]{fontenc}
\usepackage[margin=2.4cm]{geometry}
\usepackage{graphicx}\usepackage{amsmath}\usepackage{amssymb}\usepackage{xcolor}
\usepackage{booktabs}\usepackage{longtable}
\usepackage{float}\usepackage[hidelinks]{hyperref}
\usepackage{caption}\captionsetup{font=small,labelfont=bf}
\usepackage[most]{tcolorbox}

\definecolor{hypbg}{HTML}{EEF3FA}\definecolor{hyprule}{HTML}{3A6EA5}
\newtcolorbox{hyp}[1]{colback=hypbg, colframe=hyprule, boxrule=0pt, leftrule=3pt,
  arc=1pt, left=8pt, right=8pt, top=5pt, bottom=5pt, fontupper=\small,
  title={\bfseries #1}, coltitle=hyprule, fonttitle=\small\bfseries,
  attach title to upper=\par, before skip=9pt, after skip=9pt}

\definecolor{cavbg}{HTML}{FDF3E7}\definecolor{cavrule}{HTML}{C87A2F}
\newtcolorbox{caveat}{colback=cavbg, colframe=cavrule, boxrule=0pt, leftrule=3pt,
  arc=1pt, left=8pt, right=8pt, top=5pt, bottom=5pt, fontupper=\small,
  before skip=9pt, after skip=9pt}

\newcommand{\sinit}{\sigma_{\mathrm{init}}}
\newcommand{\sinput}{\sigma_{\mathrm{input}}}

\title{HIKET --- next session\\
  \large The level/trend discrepancy: what it is, and the hypotheses for it}
\author{}
\date{13 August 2026 --- working document, not for circulation}

\begin{document}
\maketitle

\begin{abstract}
\noindent Companion to \texttt{M\&M\_parameterization\_working\_document.pdf}. That document
states the parameterisation assumptions and their evidence. This one states the problem that
survives them. After correcting the error model and the pre-run anchor, the six models can
reproduce the observed soil carbon \emph{stock} but not its \emph{rate of change}: the level and
the trend make opposing demands, and no setting of the initialisation satisfies both. Six
candidate explanations are recorded here with their status and the test that would decide each: one
was implemented and rejected, two have since closed, and a seventh issue ---
the observed rate itself was not uniquely defined --- has been resolved and is retained because every
other gap is measured in its units. A final section records a \emph{second} discrepancy, the short
calibrated residence times, which turns out to be three different findings rather than one.
\end{abstract}

% =====================================================================
\section{Where the runs stand}

\textbf{Last completed run: Roihu jobs 597028--597033} (RUN\_IDs \texttt{20260812\_0809*}), testing
the error model alone --- scale $0.72\rightarrow0.80$ plus Student-$t$ tails ($\nu=6$), both measured
rather than tuned, with campaign-specific $\sigma$ deliberately excluded so it could not confound the
trend. Five of six completed; SP1 hit an out-of-memory condition at a chain boundary and all six
SLURM scripts now request 80\,GB.

\textbf{Its verdict is that this lever is exhausted.} Residence times moved under 2\,\%, the residual
spread barely moved ($0.73$--$0.75$), the returned kurtosis was $7.2$--$7.4$ against the $t$'s own
implied $6.0$, and \textbf{the 2006--2024 source survived}. The error model was worth doing --- it is
now measured and self-consistent --- but it is not the explanation for either the trend deficit or
the short residence times.

\textbf{Update 2026-08-14 --- the corrected-target run exists and is half landed.} Jobs
619204--619209 ran on the corrected target (verified in the launch log: \texttt{SOC obs: 1205},
456 calibration-ready, \texttt{SIGMA\_1985\_INFL 2.00}); \texttt{Data/} had been rsynced after all.
\textbf{SP1, TP2 and TP3 completed} (R-hat $1.000$--$1.003$, ESS $4582$--$9713$) and are synced as
\texttt{20260813\_0803*}. \textbf{Yasso07/15/20 were OOM-killed at 8\,h} and were relaunched as
jobs \textbf{654390--654392} on 2026-08-14 with identical settings, so that all six will sit on one
footing. Results expected 2026-08-15.

\textbf{The out-of-memory failures are chronic and the accounting cannot diagnose them.} The
ceiling has been raised $16\rightarrow40\rightarrow80$\,GB since July and the failures recurred at
every level, with \texttt{MaxRSS} pinned at $12$--$16$\,GB throughout --- because
\texttt{JobAcctGatherFrequency} is 10\,s and cannot see a peak. All six SLURM scripts now carry
\texttt{cgroup\_memlog.sh}, which records \texttt{memory.peak} (a kernel high-water mark, immune to
sampling) and \texttt{memory.events:max} at \emph{both} the job and step cgroup. The distinction
matters: the step's \texttt{memory.max} is \texttt{UNLIMITED}, so its breach counter would read
zero forever and look like proof of innocence. On landing, \texttt{events:max}${}>0$ means the job
breached its own ceiling and \texttt{peak} sizes the correct request; \texttt{max}${}=0$ alongside
an \texttt{oom\_kill} means the kill came from outside and more memory buys nothing. Memory was
raised to 160\,GB as headroom, explicitly not as a diagnosis.

\begin{caveat}
\textbf{Expect the re-run to narrow the problem, not to solve it.} The target corrections move the
\emph{observations} toward the models over 1985--2024 ($+0.248 \rightarrow \approx +0.182$), but
\textbf{2006--2024 is untouched at $+0.107$ against a source in all six models}, and residence times
are untouched entirely. See also \S\ref{sec:mrtgap}: the residence-time gap turns out to be three different
findings, only one of which is decisively at odds with the data.
\end{caveat}

% =====================================================================
\section{The discrepancy}

Sweeping the 1917 anchor at posterior-median parameters, the level and the trend cannot be
satisfied together:

\begin{table}[H]\centering\small
\caption{$\sinit$ sweep. Observed: 1985 $=61.6$, 2024 $=73.8$, trend $+0.312$\,tC\,ha$^{-1}$\,yr$^{-1}$.}
\begin{tabular}{lrr|rr}
\toprule
 & \multicolumn{2}{c|}{TP2} & \multicolumn{2}{c}{Yasso20}\\
$\sinit$ & 1985 & trend & 1985 & trend\\
\midrule
0.25 & 26.8 & $+0.299$ & 23.3 & $+0.194$\\
0.70 & \textbf{60.3} & $+0.137$ & 42.4 & $+0.082$\\
1.10 & 90.2 & $-0.008$ & \textbf{59.4} & $-0.018$\\
\bottomrule
\end{tabular}
\end{table}

\noindent Get the level right and the trend is a third to a tenth of observed; get the trend right
and the 1985 stock is less than half of observed. This is the problem to explain.

% =====================================================================
\section{H6 first: the target is not yet defined}

Every gap below is quoted against an observed rate that varies with two choices not yet made
explicitly.

\begin{table}[H]\centering\small
\caption{Observed paired rate of SOC stock change (tC\,ha$^{-1}$\,yr$^{-1}$), plot set defined by
the data alone.}
\begin{tabular}{lrrrr}
\toprule
period & $n$ & mean & 10\,\% trimmed & median\\
\midrule
1985--2006 & 348 & $+0.458$ & $+0.409$ & $+0.410$\\
2006--2024 & 409 & $+0.209$ & $+0.105$ & $+0.071$\\
1985--2024 & 316 & $+0.312$ & $+0.260$ & $+0.229$\\
\bottomrule
\end{tabular}
\end{table}

\begin{hyp}{H6 --- the observed rate is ambiguous by a factor of $\sim$3}
\emph{Statistic:} 1985--2024 spans $+0.229$ (median) to $+0.312$ (mean), a 36\,\% range; 2006--2024
spans a factor of three. Per-plot changes are strongly right-skewed.

\emph{Plot set:} a model-dependent subset (plots where the forward run succeeded) gave $+0.116$
for 2006--2024 where the data-defined set gives $+0.209$. A target must not depend on which plots
a model managed to run.

\smallskip
\emph{Test / action:} fix one plot set and one statistic, applied identically to observed and
modelled rates. For an inventory the \textbf{mean} is defensible --- total stock change is what is
reported, and the median discards the high-accumulation plots carrying the sink --- but the skew
should be reported alongside: a mean--median gap this large is itself a result about the
heterogeneity of accumulation.

\smallskip
\emph{Why first:} the target spans $\approx +0.07$ to $+0.31$, the same magnitude as the deficit
being explained. Some of the discrepancy may be a choice rather than a mechanism.
\end{hyp}

\begin{caveat}
\textbf{H6 --- RESOLVED 2026-08-12.} The action above has been carried out.
\texttt{doublechecks/observed\_soc\_basis.R} is now the canonical table. There are \emph{three}
axes, not one, and together they span a factor of four for the same data:
\textbf{weighting} (unweighted = the average plot in our sample; region-weighted = the average
hectare of Finland, worth 25--30\,\%), \textbf{plot set} (paired vs balanced --- for 2006--2024,
also requiring a 1985 observation moves $+0.216 \to +0.143$, a 34\,\% swing, though for 1985--2024
it changes almost nothing), and \textbf{depth} (measured 0--40\,cm vs whole profile including the
modelled deep tail, worth $\sim$35\,\% here because the tail is refitted per campaign).

\smallskip
\textbf{The rule adopted:} model comparison $\to$ \emph{paired, unweighted, whole profile};
any national statement $\to$ \emph{weighted, measured 0--40\,cm}; comparing intervals with each
other $\to$ \emph{balanced}, so the population is constant. Never difference means over different
plot subsets.

\smallskip
The table above is on the paired, unweighted, whole-profile basis and remains correct for model
comparison. \textbf{Where 2006--2024 is set beside 1985--2024, the balanced figure $+0.143$ should
be used instead of $+0.209$}, since only then do the two intervals describe the same plots.

\smallskip
\textbf{Applied to the figures the same day:} F2, F3 and F4 drew a model curve averaged over all
plots against observed means taken over per-campaign subsets. They now share one balanced set
($n=316$) via \texttt{manuscript/figures/obs\_basis.R}. On the companion appendix figure the same
mismatch had moved Yasso15's 1985 level $63.3 \to 65.1$ and its 1985--2024 rate
$+0.422 \to +0.379$, inflating the apparent model--observation gap by $\sim$40\,\%. F5 was already
correct (paired per-plot differences).
\end{caveat}

% =====================================================================
\section{The hypotheses}

\begin{hyp}{H1 --- kinetics too slow to track the rising litter}
The cascade models have an effective response time $\tau \approx 195$\,yr and close only 18\,\% of
any gap over 39 years; the observations imply $\tau \approx 33$--$38$. Yasso is fast enough
(17--38\,yr) but its rates are fixed by design, and its free transfer fractions only redistribute
carbon among pools that all turn over within a decade --- the one slow reservoir (H, $\tau\approx855$\,yr)
has a \emph{fixed} share. So Yasso's response time is close to a design constant.

\emph{Status:} identified. Freeing the slow-rate prior (SD $0.15\rightarrow0.50$) made it
\emph{worse} --- $\alpha_H$ drifted down $0.0053\rightarrow0.0031$ and the trend moved only
$0.116\rightarrow0.130$ --- because the likelihood does not reward the trend (below).

\emph{Test:} make the likelihood target the stock change, then see where the rates go. A perfect
trend fit is currently worth $3.4$ nats against a parameter excursion costing $8.9$; adding the
paired stock change as an explicit observation with its own SE ($0.035$) would be worth
$\approx13$, which flips the balance.
\end{hyp}

\begin{hyp}{H2 --- the litter trend is understated}
$\sinput$ absorbs any error in the litter \emph{level}, so only an understated \emph{rate of
increase} causes under-accumulation. The product rises $\times1.35$ over 1986--2024 against a
growing stock rising $\times1.42$.

\emph{Status:} \textbf{untested}, and the only remaining input-side hypothesis that is testable
with data in hand.

\emph{Test:} sweep a multiplier on the litter trend (holding the level) and find what rise would
be required; check it against the growing-stock rise and the fitted elasticity. Same shape as the
$\sinit$ sweep and equally decisive.
\end{hyp}

\begin{caveat}
\textbf{H2 --- CLOSED 2026-08-10, confirmed with B.~Tupek.} The series stands as published; the 2006
peak is real and not an artefact. Litter rises $+51.7$\,\% over 1986--2006 against a growing stock
rising $+22.4$\,\%, then falls $-12.8$\,\% over 2006--2021 against a growing stock rising
$+16.0$\,\% (net $\times1.32$ against $\times1.42$). \textbf{A modest post-2006 decline in modelled
SOC is therefore expected and defensible, not something to be removed.} Do not re-open. The one
remaining input-side lead is H4 (understorey), which needs external data.
\end{caveat}

\begin{hyp}{H3 --- the 1985 campaign reads low}
VMI8 samples 0--5 and 5--20\,cm where the later campaigns sample 0--10 and 10--20\,cm, and it
could not be validated against official LUKE figures. If it reads low, the observed trend is
overstated.

\emph{Status:} plausible and \textbf{not resolvable with these three campaigns.} The obvious test
--- compare over 2006--2024, which excludes 1985 --- has almost no power: the total change there is
$+3.8$\,tC\,ha$^{-1}$ on a stock of $\sim$70 ($3.2\sigma$), so a model capturing 50\,\% of the
trend is statistically indistinguishable from one capturing 100\,\%. Essentially all the
discriminating information sits in the 1985--2006 leg ($+9.6$\,tC\,ha$^{-1}$, $9.5\sigma$) --- the
contested one.

\emph{Consequence:} H1 and H3 predict the same observable. Record as a limit of the data, and as a
stated sensitivity, rather than as an open analysis task.

\emph{Exception:} Yasso20 predicts $-0.147$ against $+0.209$ over the trusted campaigns, a gap of
$5.5$\,SE. Its failure stands regardless of 1985.
\end{hyp}

\begin{caveat}
\textbf{H3 --- SUPERSEDED 2026-08-12: the defect was identified, and it is not the one hypothesised.}
Two concrete comparability failures were established from the source workbook's own structure, and
both are now corrected in the target.

\smallskip
\emph{(i) The 1985 organic layer is OFH only.} Litter-and-moss is a separately coded layer that was
measured in 2006 and 2024 and folded into their organic layer, but is \textbf{absent from 1985}. It
is worth $3.37$--$4.27$\,Mg\,C\,ha$^{-1}$ and supplies \textbf{89\,\% of the apparent 1985--2006
organic gain}. Corrected by adding each plot's own 2006 litter layer to its 1985 organic layer.

\smallskip
\emph{(ii) ``1985'' is really 1986--1995.} $19.5$\,\% of the plots were sampled in 1995; the campaign
mean represents $\approx$1989, so the true mean interval to 2024 is $34.7$ years, not $39$ --- a
12\,\% error in every rate denominator. Corrected by carrying the true sampling year per plot, and
dropping the 82 plot-years that carry no date.

\smallskip
\textbf{What is \emph{not} wrong with 1985}, against intuition: its depth distribution, its
stoniness, its organic/mineral boundary and its noise level are all unremarkable, and proportional
bias appears in \emph{every} campaign pair. Whatever remains is a \emph{level}, not scatter. A
pre-registered inflation $f=2$ on the 1985 likelihood covers the unattributable residual; no subsoil
reconstruction and no campaign bias term were adopted, the latter because a free campaign offset
consumes the very trend being measured.
\end{caveat}

\begin{hyp}{H4 --- a missing input component with its own trend}
Understorey litter is absent from the product entirely, and its share plausibly changed with
canopy closure over the period. A missing component with a rising trend would look exactly like
H2.

\emph{Status:} untested; blocked on the LUKE MUSTIKKA ground-vegetation data, which is not held
locally. Relevant that two models put $\sinput$ \emph{below} 1, against a C4b prior centred at
1.30 precisely because understorey is missing.

\emph{Test:} would appear as a productivity-dependent trend deficit, since understorey is
proportionally largest where tree litter is smallest.
\end{hyp}

\begin{hyp}{H5 --- a missing process: stabilisation capacity}
Clay adds $+0.050$ to the residual $R^2$ beyond litter ($n=456$, $p<10^{-7}$), with the right
sign: less over-prediction on higher-clay soils. No model in the set has any edaphic control on
decomposition --- a shared blind spot.

\emph{Status:} evidence in hand, mechanism not implemented. Realistically the second paper.

\emph{Caveat:} CEC adds more ($+0.087$) but is measured on the same samples as SOC and organic
matter drives CEC --- circularity. Use clay.
\end{hyp}

\begin{hyp}{H0 --- REJECTED: historical depletion of the 1917 soil}
That early-20th-century management (litter raking, forest grazing, slash-and-burn) removed organic
matter before it entered the soil, leaving 1917 stocks below what a biomass proxy implies.

\emph{Status:} \textbf{implemented and rejected} (\S2). The anchor sweep shows the level and the
trend make opposing demands: no value of $\sinit$ satisfies both, in either family. The hypothesis
may remain true as history, but it is not the mechanism behind the trend deficit.

\emph{What survives:} initialisation controls the sink \emph{sign} --- Yasso20 spans $-0.09$ to
$+0.23$ across the sweep --- but not its magnitude.
\end{hyp}

% =====================================================================
\section{The residence-time gap is three findings, not one}\label{sec:mrtgap}

Alongside the trend deficit sits a second discrepancy: our calibrated intrinsic mean residence times
come out below the published parameterisations. Work of 2026-08-13 shows this is \textbf{not one
phenomenon}, and that treating it as one has been obscuring the result.

\subsection*{What a longer residence time would cost}

In the Yasso family MRT and $\sinput$ trade off along a clear ridge ($r=-0.79/-0.73/-0.57$): the
observed stock pins their \emph{product}, so the split between ``slow turnover, modest inputs'' and
``fast turnover, inflated inputs'' must be decided by something other than the stock. The simple
models show no such ridge ($r\approx-0.26/-0.31$), so this is a property of the Yasso structure, not
of the calibration in general. Reading the best fit attainable near each published value against the
best attainable anywhere:

\begin{table}[H]\centering\small
\caption{Cost of the published residence time, in log-likelihood units, run \texttt{20260812\_0809*}.}
\begin{tabular}{lrrr}
\toprule
model & P(MRT $>$ published) & cost & draws in band\\
\midrule
Yasso07 & \textbf{0\,\%} & \textbf{never reached} & 0\\
Yasso15 & 0.011\,\% & $\le 14.7$ & 741\\
Yasso20 & \textbf{18.1\,\%} & $\le 1.6$ & 121\,141\\
\bottomrule
\end{tabular}
\end{table}

\noindent Yasso20's published value is nearly free and Yasso15 reaches its own only in the far tail.
\textbf{Yasso07 never reaches its published value at all}: 225\,015 draws span $10.4$--$25.4$\,yr
against a published $33.5$, so no sampled estimate of the cost exists --- only a profile likelihood
could supply one.

\begin{caveat}
\textbf{Two qualifications, both load-bearing.} \emph{(i)} These are maxima over posterior draws, so
they \emph{over}state the cost, and they do so most where draws are thinnest --- 121\,141 draws in
Yasso20's band, 741 in Yasso15's, \textbf{none} in Yasso07's. Settling it needs a profile likelihood
(an optimisation over the remaining parameters at each grid node, $\approx$7\,h per model), and
\textbf{Yasso07 is where to spend it}, since for that model nothing can be read off the samples.
An earlier version of this table was built on \emph{burn-in artefacts} --- the first retained iteration
of each internal chain, 15 rows per model, which were the only draws near the published value for
Yasso07 and Yasso15; they are now removed. \emph{(ii)} \textbf{Two
different ``published MRT'' values exist and must not be mixed}: the published \emph{point}
parameterisation and the median of the published posterior \emph{sample}. Residence time is a
nonlinear many-to-one function of the parameters, so these differ --- for Yasso20, $19.0$ against
$25.0$\,yr, which moves its cost from $1.6$ to $10.0$. The table uses the point. \textbf{Which
comparator the paper adopts is still to be decided.}
\end{caveat}

\subsection*{Why the climate response owns Yasso07's gap and cannot own the others'}

Yasso07 applies a \emph{single} climate modifier $\xi$ to every pool, so $\xi$ is a pure rescaling of
time and $\mathrm{MRT}=\mathrm{MRT}_{\mathrm{ref}}/\xi$ exactly. Calibration moves $\xi$ from $0.855$
to $1.822$ ($\times2.13$) at the Finnish mean climate, and that alone predicts $15.7$\,yr against an
actual $15.2$ --- the entire gap, with nothing left for the transfer fractions. The model has nowhere
else to put a climate misfit.

Yasso15 and Yasso20 carry \emph{three} pool-specific modifiers. Residence time is governed by the
slow humus pool, which has its own $\beta_{H1}$, and that modifier moved only $\times1.04$ and
$\times1.06$; in Yasso20 the components move in opposite directions ($\xi_{\mathrm{AWE}}\times0.76$
against $\xi_{\mathrm{N}}\times1.16$) and largely cancel. \textbf{The leverage Yasso07 hands to
climate is structurally unavailable to its successors.}

Two consequences. Yasso07's \emph{published} $\xi$ is $0.855<1$ --- its own parameterisation says
Finnish conditions decompose \emph{more slowly} than its reference, where Yasso15 and Yasso20 already
say faster ($1.16$--$1.66$); it starts furthest from what the data want and is the only variant able
to travel the distance in one parameter. And calibration \textbf{inverts the family ordering}:
published $33.5/30.4/19.0$\,yr (Yasso07 slowest) becomes $15.2/22.1/17.5$ (Yasso07 fastest).

\emph{Open:} whether $\xi=1.82$ is defensible. It requires $\beta_1$ to move $0.0987\rightarrow0.1578$
($+60$\,\%), some $4.4$ prior standard deviations off centre. Suggestive but not decisive: that value
is essentially Yasso20's \emph{published} $\beta_1$ ($0.1580$), though Yasso07's $\beta_1$ scales all
pools while Yasso20's scales one, so they are not strictly the same quantity. Pair with an
independent check of the fitted Finnish warming rate.

\smallskip
\noindent Reproduce with \texttt{doublechecks/xi\_published\_vs\_ours.R} and
\texttt{manuscript/figures/build\_F14\_mrt\_ridge.R}.

% =====================================================================
\section{The input-flux window is anchored on the wrong statistic}\label{sec:fluxwindow}

\emph{Established 2026-08-14. Decided but deliberately not implemented --- see the sequencing note
at the end. The full write-up, with tables and the boxed proposal, is in the M\&M working document,
section ``Prior specification: the litter-input flux window''.}

$\sinput$ is a single global scalar, so $\sinput \times \bar{J}$ is a \textbf{national mean}
effective litter flux. The current \texttt{flux\_pair} ceiling of $8.7$\,tC\,ha$^{-1}$\,yr$^{-1}$ is
Gower et al.\ (2001)'s Class~I evergreen \textbf{maximum} --- the most productive single stand in a
global boreal compilation. \textbf{The window has been bounding a mean with a maximum}, which is why
it has never constrained anything: posterior effective fluxes are $4.99$--$6.56$, all comfortably
inside. The \emph{direction} of the bound is sound --- $\text{litter}\le\text{NPP}$ is an identity
--- but the statistic and the population are wrong.

\begin{caveat}
\textbf{A unit trap worth recording, because it produced a wrong claim mid-analysis.} Gower et al.\
(2001) reports in \textbf{g\,C}\,m$^{-2}$\,yr$^{-1}$; Zheng et al.\ (2004) reports in \textbf{dry
matter}. The proof is internal to the two papers: Zheng cites Gower's world-boreal total NPP as
$109$--$1827$ (mean $892$) where Gower's own Appendix~A gives $218$--$912$ (mean $424$) g\,C --- a
ratio of exactly $2.00$--$2.10$. Zheng's headline $563$ is therefore $\approx 2.81$, not $5.63$,
tC\,ha$^{-1}$\,yr$^{-1}$. Mixing the two inflates any NPP ceiling twofold.
\end{caveat}

\begin{center}\small
\begin{tabular}{@{}lrll@{}}
\toprule
window ceiling & $\sinput\le$ & binds? & implied MRT Y07/Y15/Y20 \\
\midrule
$8.70$ current (Gower global max) & 3.46 & \textbf{no} & 15.2 / 21.9 / 17.5 (unchanged) \\
$4.62$ Gower Nordic max \quad$\leftarrow$ \textbf{decided} & \textbf{1.84} & all six & 21.3 / 25.8 / 24.0 \\
$2.81$ Zheng gridded mean & 1.12 & all six & 35.0 / 42.4 / 39.5 \\
\midrule
\multicolumn{4}{@{}l@{}}{\emph{published, for reference:} 33.5 / 30.4 / 19.0} \\
\bottomrule
\end{tabular}
\end{center}

A constraint \emph{binds} when the value the likelihood would choose lies outside it, so that the
bound rather than the data sets the answer. Zheng's ceiling is not adoptable as a drop-in: it lies
\emph{below} the prior centre ($\sinput = 1.30 \Rightarrow F = 3.26$), so the centre is not
representable; and being only $4\%$ above LUKE's own total litter, it amounts to asserting the
inventory as near-exact. Narrowing the window also \textbf{tightens the prior}, not only the wall:
at fixed \texttt{sigma\_ppm} the $\pm1$\,ppm range moves from $\sinput\in[0.93,1.72]$ to
$[1.09,1.47]$, because the transform is a scaled logit onto the window.

\begin{caveat}
\textbf{Three claims that did not survive scrutiny.} (i) ``our flux is physically impossible'' --- it
exceeds the Nordic \emph{mean} but not the Nordic \emph{range} ($2.15$--$4.62$); (ii) ``LUKE's litter
product is biased high'' --- $2.70$ is $84\%$ of Gower's Nordic mean, but that range is too wide to
support the inference; (iii) ``Finnish forests are twice as productive as Gower's stands'' --- this
compared Korhonen's \emph{current annual increment} against Gower's \emph{mean annual increment}
(biomass $\div$ stand age, on stands averaging 99 years), so the factor is largely definitional, and
Gower's stands are selected for having complete NPP budgets rather than for representing Finnish
managed forest. A national inventory rebuilt across thirteen cycles is not likely to be wrong by a
factor of two without anyone noticing; the discrepancy was in the comparison.
\end{caveat}

\textbf{What $\bar{J}$ contains, now established} from the Zenodo deposit (DOI
10.5281/zenodo.19736499, now live) and the \texttt{YaYasso} implementation of the LUKE chain: it is
a flux (biomass $\times$ turnover), it \textbf{includes harvest residues and natural mortality} ---
both previously assumed missing --- and it \textbf{excludes understorey}, which is what the
$\sinput$ centre of $1.30$ stands in for. The dominant uncertainty is the fine-root chain, where
\texttt{YaYasso} variants span foliage$\times 0.18$--$2.5$ for biomass and $\times0.5$--$0.85$ for
turnover: a $\sim$20-fold range in the dominant component, and precisely the term Lehtonen \&
Heikkinen exclude from their own uncertainty, which makes their $\pm10\%$ a floor.

\textbf{Sequencing.} Wait for jobs 654390--654392, analyse them, \emph{then} decide. Do not bundle
this with the $\sinput$ prior tightening or anything else --- one change per run is the discipline
that produced every clear result so far.

% =====================================================================
\section{Suggested order}

\emph{Revised 2026-08-13.} H6, H2 and H3 have all closed since the list was first written; what
remains is reordered accordingly.

\begin{enumerate}
\item \textbf{Re-run the six calibrations on the corrected target.} Everything downstream is quoted
      against the pre-correction observations, so no other number can be trusted to move for the
      right reason until this lands. Push and rsync first.
\item \textbf{Decide the published-MRT comparator} (point versus posterior median). Free, and it
      changes a headline number by a factor of seven for Yasso20.
\item \textbf{Profile likelihood for Yasso07} --- the only way to price a longer residence time
      without the prior contaminating the coverage, and the only way to tell whether the
      $\le137.7$ figure is real or an artefact of thin sampling.
\item \textbf{H1} --- change what the likelihood targets, then re-read the rates. Expensive, and its
      answer depends on (1). Note \S\ref{sec:mrtgap}: for Yasso07 the rates are not where the gap lives, so this
      test now applies mainly to the simple models and to Yasso15.
\item \textbf{H4, H5} --- second paper, or as discussion; H4 needs external data.
\end{enumerate}

\begin{caveat}
These are not exclusive. Several could contribute simultaneously, and the ordering above is by cost
and dependency, not by expected importance. In particular H1 and H3 are observationally degenerate
over the available campaigns, so neither can be confirmed by ruling the other out --- and the H3
corrections, while they remove a real defect, do not resolve that degeneracy.

\smallskip
\textbf{The load-bearing caution now.} The error-model lever is exhausted, the input series is
confirmed, and the target is corrected --- yet 2006--2024 remains a source in all six models against
an observed sink. The remaining candidates are structural, and \S\ref{sec:mrtgap} shows they differ by model. Resist
any single explanation for what are, on the evidence, several different problems.
\end{caveat}

\end{document}
